\section{Related Work}
\label{sec:related_work}

\textbf{Model Merging and Parameter Fusion.} Weight-space model merging has gained significant traction as a resource-efficient alternative to multi-task learning, allowing the combination of multiple downstream models into a single base architecture without joint training \cite{ilharco2022editing, wortsman2022model}. Early methods explored linear interpolation in weight space, such as Model Soups \cite{wortsman2022model} or Task Arithmetic \cite{ilharco2022editing}. To address structural misalignments between independent training runs, permutation-based matching techniques like Git Re-Basin \cite{ainsworth2022git} and ZipIt! \cite{zhang2023zipit} align the neurons of disjoint networks before fusion. 

To combat destructive parameter interference, TIES-Merging \cite{yadav2023ties} filters out low-magnitude parameter updates and resolves sign conflicts before averaging. More recently, adaptive model merging frameworks like AdaMerging \cite{yang2023adamerging} and RegCalMerge \cite{regcalmerge2024} learn layer-wise linear coefficients on small calibration datasets using unsupervised prediction objectives. However, because they optimize unconstrained unregularized layer-wise coefficients directly in real parameter space, they are prone to the Overfitting-Optimizer Paradox under extreme data scarcity. While PolyMerge \cite{polymerge2024} and SplineMerge \cite{splinemerge2024} address this by constraining the search space to low-dimensional depth-wise subspaces, they remain restricted to real-valued Euclidean updates. PhaseMerge completely rejects real-space scalar constraints, operating instead in the complex-valued Fourier domain.

\textbf{Post-Training Quantization (PTQ) and Q-Merge.} Model compression is vital for deploying heavy transformer models on resource-constrained hardware \cite{gholami2021survey}. Post-training quantization (PTQ) enables low-bit network representations without the high overhead of quantization-aware training \cite{frantar2022gptq, lin2023awq}. Standard techniques, including GPTQ \cite{frantar2022gptq}, AWQ \cite{lin2023awq}, SmoothQuant \cite{xiao2023smoothquant}, and LLM.int8() \cite{dettmers2022llm}, optimize scale and rounding parameters under low-bit schemas. However, the discrete nature of rounding operators introduces highly non-linear, discontinuous landscapes that are notoriously difficult to navigate. 

To merge models under quantization, Q-Merge \cite{qmerge2024} optimizes merging coefficients directly under PTQ constraints. However, because Q-Merge's learned parameters are highly sensitive to specific rounding boundaries of the calibration schema, they exhibit poor generalization under target deployment schema shifts. By contrast, PhaseMerge's low-dimensional phase-shift grids generate smooth parameter perturbations in frequency space that are inherently robust to rounding boundaries, ensuring excellent generalizability across diverse target schemas.

\textbf{Fourier and Spectral Methods in Deep Learning.} Mathematical transforms in the frequency domain have long been utilized to compress and analyze deep neural networks. Rippel et al.\ \cite{rippel2015spectral} demonstrated that convolutional networks can be parameterized in the spectral domain using Fast Fourier Transforms (FFT) for faster convergence and parameter compression. Similarly, Fourier Neural Operators (FNOs) \cite{li2020fourier} and Fast Fourier Convolutions \cite{chi2020fast} model continuous physical systems and long-range spatial dependencies in the frequency domain. Frequency-space representations have also been crucial in coordinate networks, such as NeRF \cite{mildenhall2020nerf, tancik2020fourier}, where high-frequency positional embeddings enable networks to learn fine-grained spatial details.

Closest to our work, FREE-Merging \cite{freemerge2024} applies static, non-adaptive Fourier high-pass and low-pass filters to task vectors to filter out redundant parameters and reduce inter-task interference. However, FREE-Merging is non-adaptive, relying on hard-coded frequency boundaries that do not adjust to task dynamics. PhaseMerge is fundamentally different: instead of static frequency-space filtering, we introduce \emph{differentiable, continuous phase-rotation grids} that are optimized end-to-end. This is the first model-merging framework that dynamically learns constructive and destructive interference patterns to actively cancel out parameter conflicts and quantization noise.
